% =============================================================================
% The LQ Digital Mode Family — Reference Implementation (lq_lib)
%
% Author:  Luis Quesada (HB9IPH)
% Web:     https://luisquesada.com
% Portal:  https://lquesada.github.io/lq_lib/
% GitHub:  https://github.com/lquesada/lq_lib
% App:     qFT8 — Portable Amateur Radio for Android (https://qft8.com)
%
% License: MIT License (https://github.com/lquesada/lq_lib/blob/main/LICENSE)
%
% Copyright (c) 2026 Luis Quesada (HB9IPH)
%
% Permission is hereby granted, free of charge, to any person obtaining a copy
% of this software and associated documentation files (the "Software"), to deal
% in the Software without restriction, including without limitation the rights
% to use, copy, modify, merge, publish, distribute, sublicense, and/or sell
% copies of the Software, and to permit persons to whom the Software is
% furnished to do so, subject to the following conditions:
%
% The above copyright notice and this permission notice shall be included in
% all copies or substantial portions of the Software.
%
% THE SOFTWARE IS PROVIDED "AS IS", WITHOUT WARRANTY OF ANY KIND, EXPRESS OR
% IMPLIED, INCLUDING BUT NOT LIMITED TO THE WARRANTIES OF MERCHANTABILITY,
% FITNESS FOR A PARTICULAR PURPOSE AND NONINFRINGEMENT. IN NO EVENT SHALL THE
% AUTHORS OR COPYRIGHT HOLDERS BE LIABLE FOR ANY CLAIM, DAMAGES OR OTHER
% LIABILITY, WHETHER IN AN ACTION OF CONTRACT, TORT OR OTHERWISE, ARISING FROM,
% OUT OF OR IN CONNECTION WITH THE SOFTWARE OR THE USE OR OTHER DEALINGS IN THE
% SOFTWARE.
% =============================================================================

% ===========================================================================
%  Section VII: Conclusion, Acknowledgment & References (ALLOCATION: PAGE 8)
%  Must fit entirely on Page 8:
%    - Table 7 (Multi-Mode Comparison) [Top]
%    - Left Column: Section VII (Conclusion) & Acknowledgment
%    - Right Column: The Bibliography References [1]–[13]
% ===========================================================================

\section{Conclusion}
\label{sec:conclusion}

\textbf{LQ8} demonstrates that continuous-phase 8-GFSK and systematic LDPC$(174,91)$ transport deliver higher throughput through variable-length prefix framing while maintaining full backward physical compatibility and $-21.0$\,dB SNR sensitivity in a 2500\,Hz reference bandwidth. By optimizing message prefix allocations to match operational entropy, \LQ8 packs callsigns, locators, and reports into a 77-bit budget, finalizing confirmed two-way contacts in 4 transmission slots (60\,s). For high-density pileups, \LQ8 introduces the single-carrier \msgtype{MULTI-REPORT+73} protocol, confirming two answering stations simultaneously without RF power-splitting penalties and achieving throughputs up to 240\,QSOs/h (480\,QSOs/h in dual-carrier operation).

The unified protocol family extends this architecture across three complementary physical profiles. \LQ16 scales symbol duration to 320\,ms, delivering a $+3.0$\,dB sensitivity advantage ($-24.0$\,dB SNR threshold) across ultra-weak fading and EME paths. For rapid contesting and dynamic propagation, \LQ4 (7.5\,s slots) and \LQ2 (3.75\,s slots) utilize 4-GFSK modulation and Costas sync arrays to complete contacts in 30.0\,s and 15.0\,s, supporting burst rates up to 960\,QSOs/h on VHF/UHF meteor scatter and sporadic-E openings.

The complete software suite is implemented in the open-source C++ reference library (\texttt{lq\_lib}) under the MIT License, featuring an automated verification harness achieving $>90\%$ line coverage alongside comprehensive golden vector test suites across all operational profiles. End-to-end field verification has been demonstrated in real-world portable activations via \textbf{qFT8} for Android across HF bands. Future protocol evolution will explore multi-carrier wideband chat aggregation, dedicated 60-meter channel sub-allocations, and low-complexity satellite telemetry framing.

\section*{Acknowledgment}
The original concept, protocol architecture, and reference implementations were developed by the author. Generative AI was used to assist with background research, code testing, and manuscript drafting. The author thanks the global amateur radio community for valuable feedback during early on-air testing.

\pagebreak

\begin{thebibliography}{99}
\fontsize{7.6pt}{8.8pt}\selectfont
\raggedright
\interlinepenalty=10000
\setlength{\itemsep}{2.2pt plus 0.6pt minus 0.4pt}
\setlength{\parskip}{0pt}
\setlength{\parsep}{0pt}
\bibitem{ft8}
  S.~Franke, K9AN, B.~Somerville, G4WJS, and J.~Taylor, K1JT,
  ``The FT4 and FT8 Communication Protocols,''
  \textit{QEX}, pp.~7--17, Jul./Aug.\ 2020.

\bibitem{wsjtx}
  J.~Taylor, K1JT, S.~Franke, K9AN, and B.~Somerville, G4WJS,
  ``WSJT-X: Weak-Signal Communication by K1JT,''
  \mbox{\url{https://wsjtx.github.io/wsjtx/}}, 2024.

\bibitem{shannon}
  C.~E.~Shannon,
  ``A Mathematical Theory of Communication,''
  \textit{Bell System Technical Journal}, vol.~27, no.~3,
  pp.~379--423, Jul.\ 1948.

\bibitem{huffman}
  D.~A.~Huffman,
  ``A Method for the Construction of Minimum-Redundancy Codes,''
  \textit{Proceedings of the IRE}, vol.~40, no.~9,
  pp.~1098--1101, Sep.\ 1952.

\bibitem{gallager}
  R.~G.~Gallager,
  ``Low-density parity-check codes,''
  \textit{IRE Transactions on Information Theory}, vol.~8, no.~1,
  pp.~21--28, Jan.\ 1962.

\bibitem{costas}
  J.~P.~Costas,
  ``A study of a class of detection waveforms having nearly ideal range-Doppler ambiguity properties,''
  \textit{Proceedings of the IEEE}, vol.~72, no.~8,
  pp.~996--1009, Aug.\ 1984.

\bibitem{varicode}
  P.~Martinez, G3PLX,
  ``PSK31: A new digital mode for HF and VHF amateur radio communications,''
  \textit{RadCom}, vol.~75, no.~1, pp.~42--46, Jan.\ 1999.

\bibitem{maidenhead}
  IARU Region~1 (J.~Morris, GM4ANB, ed.),
  ``Maidenhead Locator System,''
  \mbox{\url{https://www.iaru-r1.org/}}, 1999.

\bibitem{crc24q}
  RTCA, Inc.,
  ``Minimum Operational Performance Standards for 1090 MHz Extended Squitter Automatic Dependent Surveillance-Broadcast (ADS-B),''
  RTCA DO-260B / ICAO Doc 9871, Dec.\ 2009.

\bibitem{lewand}
  R.~E.~Lewand,
  \textit{Cryptological Mathematics},
  Mathematical Association of America, 2000.

\bibitem{lqlib}
  L.~Quesada, HB9IPH,
  ``LQ Digital Mode Family --- C++ Reference Library (lq\_lib),''
  \mbox{\url{https://github.com/lquesada/lq_lib}}, 2026.

\bibitem{lqportal}
  L.~Quesada, HB9IPH,
  ``The LQ Digital Mode Family Portal,''
  \mbox{\url{https://lq8.org}}, 2026.

\bibitem{qft8}
  L.~Quesada, HB9IPH,
  ``qFT8: Portable Amateur Radio for Android,''
  \mbox{\url{https://qft8.com}}, 2026.
\end{thebibliography}
